\section{Preliminaries and related work}
\label{sec:background}
\label{sec:prerequisites}

\subsection{Simplicial sets, horns, and quasicategories}

\begin{definition}[Simplicial set]
Let $\Delta$ denote the category whose objects are the finite nonempty linear
orders $[n] = \{0 < 1 < \dots < n\}$ and whose morphisms are order-preserving
maps. A \emph{simplicial set} is a functor $X : \Delta^{\mathrm{op}} \to
\mathsf{Set}$. The set $X_n = X([n])$ is the set of \emph{$n$-simplices};
$X_0$ are \emph{vertices}, $X_1$ are (directed) \emph{edges}. The category of
simplicial sets is denoted $\sset$.
\end{definition}

\begin{example}
The \emph{standard $n$-simplex} $\Delta^n = \Delta(-,[n])$. Concretely,
$\Delta^2$ has three vertices $0,1,2$, three edges $0\to1$, $1\to2$,
$0\to2$, and one nondegenerate $2$-simplex filling the triangle.
\end{example}

\begin{definition}[Horn]
For $0 \le k \le n$, the \emph{horn} $\Lambda^n_k \subset \Delta^n$ is the
subcomplex obtained by removing the interior of $\Delta^n$ and the face
opposite vertex $k$. A horn is \emph{inner} if $0 < k < n$.
\end{definition}

\begin{example}
$\Lambda^2_1$ consists of two composable edges $0 \to 1 \to 2$ with the
composite $0 \to 2$ and the filling triangle \emph{missing}. Filling
$\Lambda^2_1$ means supplying both.
\end{example}

\begin{definition}[Kan complex; quasicategory; nerve]
A simplicial set $X$ is a \emph{Kan complex} if every map $\Lambda^n_k \to X$
extends along $\Lambda^n_k \hookrightarrow \Delta^n$ for all horns, and a
\emph{quasicategory} if this holds for all \emph{inner} horns. The
\emph{nerve} $\N(\mathcal{C})$ of a small category $\mathcal{C}$ is the
simplicial set with $\N(\mathcal{C})_n = $ chains of $n$ composable
morphisms. A simplicial set is a nerve of a category iff every inner horn has
a \emph{unique} filler \citep{lurie2009htt}.
\end{definition}

\begin{remark}[The spectrum used in this paper]
Bare directed graphs (only $0$- and $1$-cells): nothing filled.
Quasicategories: everything filled, fillers unique only up to higher cells.
Nerves: everything filled uniquely. We will interpret these respectively as
the newborn universe, the completed universe, and the classical
(zero-information) limit.
\end{remark}

\begin{remark}[Opposites]
The involution of $\Delta$ reversing each linear order induces $X \mapsto
X^{\mathrm{op}}$, under which $\Lambda^n_k$ becomes $\Lambda^n_{n-k}$. The two
outer horns are exchanged and inner horns go to inner horns, so $X$ is a
quasicategory exactly when $X^{\mathrm{op}}$ is. Reversing orientations
therefore cannot turn an inner-horn filling policy into an outer-horn one.
\end{remark}

\subsection{Rewriting, events, and causal order}

\begin{definition}[Non-deleting rewriting rule]
A \emph{rule} is a monomorphism $L \hookrightarrow R$ of finite simplicial
sets. An \emph{application} of the rule to a configuration $X$ is a choice of
match $m : L \to X$; the result is the pushout $X \cup_L R$. Because $L
\hookrightarrow R$ is a monomorphism and nothing is removed, applications
only add cells. (This is the non-deleting special case of double-pushout
graph rewriting \citep{ehrig2006dpo}, where the interface equals the
left-hand side.)
\end{definition}

\begin{definition}[Run; event; causal order]
\label{def:causalorder}
A \emph{run} is a finite sequence of rule applications starting from an
initial object. Each application is an \emph{event}. Event $b$ \emph{depends}
on event $a$ (written $a \prec b$) if the match of $b$ uses a cell created by
$a$; the reflexive-transitive closure $\le$ of $\prec$ is the \emph{causal
order}. The \emph{causal past} of an event $e$ is
$\past{e} = \{a : a \le e\}$.
\end{definition}

\begin{remark}[Orientation is not the causal order]
Two directions must be kept apart. The \emph{orientation} of an edge is part
of the simplicial set and decides which horns are inner. The causal order of
Definition~\ref{def:causalorder} is an order on events and records which move
supplied the cells another move consumed. Since no rule deletes anything, the
causal order is a partial order whatever the orientations do, and a branch
whose edges all point towards its origin still grows away from it.
\end{remark}

\begin{remark}
The triple (events, causality, independence) is an event structure in the
sense of \citet{winskel1986event}; we use only the partial order.
\end{remark}

\begin{definition}[Linearization; frame]
A \emph{linearization} (or \emph{linear extension}) of a finite poset
$(P,\le)$ is a total order on $P$ extending $\le$. We write $\Ext(P)$ for the
number of linear extensions. In this framework a \emph{reference frame} is a
linearization of the causal order.
\end{definition}

\subsection{Associahedra and the coherence tower}

\begin{definition}[Catalan numbers; associahedron]
The $n$-th Catalan number is $C_n = \binom{2n}{n}/(n+1)$; $C_{n-1}$ counts
the full bracketings of $n$ composable edges. The \emph{associahedron} $K_n$
\citep{stasheff1963} is the polytope whose vertices are these bracketings
and whose faces encode partial bracketings; $K_4$ is the pentagon.
\end{definition}

The relevance is dynamical: filling $2$-horns creates composites, whose
alternative bracketings pose $3$-horn questions, whose answers pose
compatibility ($4$-horn, pentagon) questions, and so on. Completion promotes
questions up this tower; it never depletes them. Since the number of
bracketings grows like $C_{n-1} \sim 4^n / n^{3/2}$, question supply grows
exponentially with tower level.

\subsection{Related work}

\gap{This subsection must be revised after the adversarial literature scan is
complete. Overlap verdicts below are provisional.}

\paragraph{Discrete gravity.} Regge \citep{regge1961} showed that curvature
on a piecewise-flat complex is concentrated on hinges as deficit angles and
that the corresponding action converges to the Einstein--Hilbert action.
Dynamical triangulations freeze edge lengths and make the gluing
combinatorics dynamical; causal dynamical triangulations
\citep{ambjorn2012cdt} additionally impose a global time foliation and
recover a de Sitter-like volume profile numerically. Our framework shares
the combinatorial ontology but differs in three respects: moves are
non-deleting (CDT's Monte Carlo moves are not), no foliation is imposed
(causal order is derived from move dependency), and the object grown is a
directed simplicial set headed toward a quasicategory rather than a
triangulated pseudomanifold of fixed topology.

\paragraph{Causal sets.} Causal set theory \citep{bombelli1987} takes a
locally finite partial order as the fundamental ontology, with
order-plus-number determining geometry, and classical sequential growth
\citep{rideout2000csg} supplies a stochastic birth process. Our causal
order is likewise fundamental-but-derived (from move dependency), and our
sprouting move is a birth event; the essential difference is that our
elements are not featureless points but cells of a simplicial set, so that
composition, coherence, and the entropy of Section~\ref{sec:entropy} have
combinatorial substrate. To our knowledge the linear-extension entropy of a
causal past has not been used as a dynamical functional in the causal set
literature. \gap{verify this claim in the lit scan}

\paragraph{Rewriting cosmologies.} The Wolfram physics project
\citep{gorard2020} evolves hypergraphs by rewriting rules and argues that
\emph{causal invariance} (confluence of the rewriting system) yields
observer-independence of the causal graph and, conjecturally, Lorentz
covariance. We adopt the same event/causal-graph architecture and the same
open problem (Section~\ref{sec:open}); our divergences are the specific
simplicial move set (horn filling, with the quasicategory as limit object),
the derived rather than postulated rule (Section~\ref{sec:maxent}), and the
entropy functional.

\paragraph{Emergent and thermal time.} That time may be state-dependent
is the thermal time hypothesis of \citet{connes1994thermal}, built on
Tomita--Takesaki modular flow; that time may be entanglement is
\citet{page1983}. Our proposal is coarser but of the same family: time is
identified with an irreversible completion process, and equilibrium (the
quasicategory limit, no unfilled inner horns) is the state in which the
process halts.

\paragraph{CPT-symmetric cosmology.} Pairing a universe with its mirror image
so that the total object carries no arrow is the combinatorial counterpart of
the universe-antiuniverse pair of \citet{boyle2018cpt} and of the two arrows
emerging from a shared middle in the Janus point picture of
\citet{barbour2014janus}. In both of those the mirror branch is determined by
the branch rather than evolving independently, which is also our convention.

\paragraph{Entropic derivations.} Jaynes' maximum-entropy inference
\citep{jaynes1957}, its dynamical extension by \citet{caticha2011}, and
Jacobson's derivation of the Einstein equations from the Clausius relation
on local causal horizons \citep{jacobson1995} together define the strategy
of Sections~\ref{sec:maxent} and~\ref{sec:open}: derive the rule by
inference, and aim the long game at an area law.

\paragraph{Computation-budget kinematics.} The Margolus--Levitin bound
\citep{margolus1998} caps the rate of distinguishable state changes by
energy, and \citet{lloyd2002} applies it cosmologically. Our budget axiom
(Section~\ref{sec:kinematics}) is of this lineage.
